% Prepared for Edward Baker, 22 October 2026.
% LLM assistance: gpt-6-astra (Codex), effort ultra, inherited parent configuration.
% Standalone manuscript fragment: requires amsmath, amssymb, amsthm and
% theorem/lemma/proposition/corollary/remark environments in the main document.

\section{Finite pair matrices and the complete heat dynamics}
\label{pm:section}

A pair statistic becomes useful for this problem only after its
population, spatial normalization, singular weights, and exterior
contribution have been specified.  This section works with a fixed finite
population.  Its matrix is the inverse-square graph Laplacian, which is
also the Hessian of the logarithmic pair energy.  The subsequent canonical
theta application proves convergence to a clock matrix by a direct
complete-source concentration estimate.  No random ensemble or statistical
universality assumption enters either argument.

\subsection{A finite inverse-square matrix}
\label{pm:matrix-subsection}

For an ordered configuration
\(x=(x_1,\ldots,x_m)\), where \(x_1<\cdots<x_m\) and \(m\ge2\), set
\begin{equation}
 S_i(x)=\sum_{j\ne i}\frac1{x_i-x_j},\qquad
 W_{ij}(x)=\frac1{(x_i-x_j)^2}\quad(i\ne j).
 \label{pm:S-and-W}
\end{equation}
Define the real symmetric matrix \(\mathcal L(x)\) by
\begin{equation}
 \mathcal L(x)_{ij}=-W_{ij}(x)\quad(i\ne j),\qquad
 \mathcal L(x)_{ii}=\sum_{j\ne i}W_{ij}(x).
 \label{pm:laplacian-definition}
\end{equation}
For every \(v\in\mathbb R^m\),
\begin{equation}
 v^{\mathsf T}\mathcal L(x)v
   =\sum_{i<j}\frac{(v_i-v_j)^2}{(x_i-x_j)^2}.
 \label{pm:quadratic-form}
\end{equation}
In particular \(\mathcal L(x)\) is positive semidefinite, its kernel is
exactly the constant vectors, and all its other eigenvalues are positive.
The positivity is a statement about a configuration at one time; by
itself it gives no sign for the time derivative of this matrix.

\begin{theorem}[Finite clock-matrix comparison]
\label{pm:clock-matrix}
Let \(c_i=c_0+i-1\), \(1\le i\le m\), and let
\(\mathcal L_* =\mathcal L(c)\).  Suppose
\begin{equation}
 |y_i-c_i|\le\delta\quad(1\le i\le m),\qquad 0\le\delta<\tfrac12.
 \label{pm:clock-error}
\end{equation}
Then the points \(y_i\) are strictly increasing and
\begin{equation}
 (1+2\delta)^{-2}\mathcal L_*
   \preceq\mathcal L(y)
   \preceq(1-2\delta)^{-2}\mathcal L_*.
 \label{pm:loewner-comparison}
\end{equation}
Here \(A\preceq B\) means that \(B-A\) is positive semidefinite.
Furthermore,
\begin{equation}
 \|\mathcal L(y)-\mathcal L_*\|_{\mathrm{op}}
 \le16\delta\sum_{d=1}^{m-1}(d-2\delta)^{-3}.
 \label{pm:operator-comparison}
\end{equation}
Writing \(0=\lambda_1<\lambda_2\le\cdots\le\lambda_m\) for the
ordered eigenvalues, one has
\begin{equation}
 \lambda_2(\mathcal L_*)\ge\frac{m}{(m-1)^2},\qquad
 \|\mathcal L_*\|_{\mathrm{op}}
   \le4\sum_{d=1}^{m-1}d^{-2}.
 \label{pm:clock-spectrum}
\end{equation}
Every nonzero eigenvalue of \(\mathcal L(y)\) satisfies the corresponding
multiplicative bounds in \eqref{pm:loewner-comparison}.
\end{theorem}

\begin{proof}
If \(i<j\) and \(d=j-i\), then
\begin{equation}
 d-2\delta\le y_j-y_i\le d+2\delta.
 \label{pm:gap-comparison}
\end{equation}
Since \(d\ge1\), these inequalities imply positivity of every adjacent
gap and
\[
 (1-2\delta)d\le y_j-y_i\le(1+2\delta)d.
\]
Insert the resulting inverse-square bounds into
\eqref{pm:quadratic-form} to obtain
\eqref{pm:loewner-comparison}.

On the positive interval between \(d\) and \(y_j-y_i\), the derivative
of \(r^{-2}\) has absolute value at most \(2(d-2\delta)^{-3}\).
The mean-value theorem therefore gives
\[
 \left|W_{ij}(y)-\frac1{(j-i)^2}\right|
 \le\frac{4\delta}{(|j-i|-2\delta)^3}.
\]
The sum of the absolute values in a row of a Laplacian difference is at
most twice the sum of its absolute weight differences.  There are at most
two indices at each label distance \(d\).  Consequently every absolute
row sum is bounded by the right side of
\eqref{pm:operator-comparison}.  For a real symmetric matrix its operator
norm is at most its maximum absolute row sum, proving that estimate.

Every clock weight is at least \((m-1)^{-2}\).  If
\(\sum_i v_i=0\), the elementary identity
\[
 \sum_{i<j}(v_i-v_j)^2=m\sum_i v_i^2
\]
therefore proves the first bound in \eqref{pm:clock-spectrum} by the
Rayleigh quotient.  The diagonal degree of any clock vertex is at most
\(2\sum_{d=1}^{m-1}d^{-2}\), so its absolute Laplacian row sum is at most
twice that quantity.  This proves the second bound.  The eigenvalue
comparisons follow from the minimum--maximum characterization of the
ordered eigenvalues on the common orthogonal complement of constants.
\end{proof}

The physical normalization matters.  If \(y=\alpha x\), \(\alpha>0\),
then
\begin{equation}
 S(x)=\alpha S(y),\qquad
 \mathcal L(x)=\alpha^2\mathcal L(y).
 \label{pm:spatial-scaling}
\end{equation}
Thus convergence of \(\mathcal L(y)\) to a finite matrix is consistent
with divergence of the physical inverse-square weights if
\(\alpha\to\infty\).  The convergence statement and its spatial scale
must always be given together.

\subsection{Complete root motion and the logarithmic energy}
\label{pm:root-motion-subsection}

Let \(H(t,z)\) be continuously differentiable in real \(t\), holomorphic
in \(z\), real for real \(t,z\), and satisfy
\begin{equation}
 H_t=-H_{zz}.
 \label{pm:heat-equation}
\end{equation}
On a real time interval \(I\), let \(x_1(t)<\cdots<x_m(t)\) be distinct
simple real zeros, with their labels continued continuously.  Implicit
differentiation makes each label continuously differentiable.  Put
\[
 P(t,z)=\prod_{i=1}^m(z-x_i(t)),\qquad H(t,z)=U(t,z)P(t,z).
\]
The quotient is interpreted by removable continuation at the selected
roots.  It is nonzero near each of them.  Define
\begin{equation}
 b_i(t)=\frac{U_z(t,x_i(t))}{U(t,x_i(t))}.
 \label{pm:exterior-unit}
\end{equation}
These numbers are real.  They retain every zero outside the selected
finite population and every nonzero factor of the field.  No infinite
root sum is being truncated or reordered.  A globally proper cluster
provides one convenient way to construct \(U\), but the calculations
below require only these regular local factorizations along the roots.

\begin{proposition}[Complete finite-cluster dynamics]
\label{pm:complete-root-dynamics}
Under the preceding assumptions,
\begin{equation}
 x_i'=2S_i(x)+2b_i.
 \label{pm:root-equation}
\end{equation}
For \(i\ne j\), one has
\begin{align}
 \frac{x_i'-x_j'}{x_i-x_j}
  &=\frac4{(x_i-x_j)^2}
   -2\sum_{k\ne i,j}
        \frac1{(x_i-x_k)(x_j-x_k)}
   +2\frac{b_i-b_j}{x_i-x_j},
 \label{pm:gap-evolution}\\
 W_{ij}(x)'&=-2W_{ij}(x)\frac{x_i'-x_j'}{x_i-x_j}.
 \label{pm:physical-weight-evolution}
\end{align}
\end{proposition}

\begin{proof}
Differentiate \(H(t,x_i(t))=0\), use
\eqref{pm:heat-equation}, and divide by the nonzero \(H_z(t,x_i)\):
\[
 x_i'=\frac{H_{zz}(t,x_i)}{H_z(t,x_i)}.
\]
At a root, \(H_z=UP_z\) and \(H_{zz}=UP_{zz}+2U_zP_z\).
Writing \(P=(z-x_i)\prod_{j\ne i}(z-x_j)\) gives
\[
 \frac{P_{zz}(x_i)}{P_z(x_i)}
   =2\sum_{j\ne i}\frac1{x_i-x_j}.
\]
This proves \eqref{pm:root-equation}.  Next,
\[
 S_i-S_j=\frac2{x_i-x_j}
 -(x_i-x_j)\sum_{k\ne i,j}
       \frac1{(x_i-x_k)(x_j-x_k)}.
\]
Subtracting the two root equations proves
\eqref{pm:gap-evolution}.  Ordinary differentiation of
\((x_i-x_j)^{-2}\) proves
\eqref{pm:physical-weight-evolution}.
\end{proof}

Define the squared Vandermonde and its negative logarithm by
\begin{equation}
 \Delta(x)=\prod_{i<j}(x_j-x_i)^2,
 \qquad E(x)=-\log\Delta(x).
 \label{pm:energy-definition}
\end{equation}
These are functions on the ordered simple-root chamber.  The energy is
singular at a collision; none of its differential formulas below is
asserted at the collision itself.

\begin{lemma}[Pair identities and the forced energy balance]
\label{pm:forced-energy}
On that chamber,
\begin{equation}
 \nabla E=-2S,\qquad \nabla^2E=2\mathcal L(x),\qquad
 \|S(x)\|^2=2\sum_{i<j}(x_i-x_j)^{-2}.
 \label{pm:energy-identities}
\end{equation}
Along the complete heat motion,
\begin{align}
 E'
  &=-4\|S\|^2-4S\mathbin{\cdot}b\notag\\
  &=-8\sum_{i<j}(x_j-x_i)^{-2}
    -4\sum_{i<j}\frac{b_j-b_i}{x_j-x_i}.
 \label{pm:forced-energy-balance}
\end{align}
Equivalently,
\begin{equation}
 (\log\Delta)'=P_{\mathrm{pair}}+B_{\mathrm{ext}},\qquad
 P_{\mathrm{pair}}=8\sum_{i<j}(x_j-x_i)^{-2},\quad
 B_{\mathrm{ext}}=4\sum_{i<j}\frac{b_j-b_i}{x_j-x_i}.
 \label{pm:complete-discriminant-balance}
\end{equation}
\end{lemma}

\begin{proof}
Differentiating \(-2\sum_{i<j}\log(x_j-x_i)\) gives the gradient
and Hessian.  In the expansion of \(\sum_i S_i^2\), the diagonal
summands give twice the inverse-square pair sum.  For every unordered
triple of distinct coordinates \(a,b,c\), the remaining terms cancel by
\[
 \frac1{(a-b)(a-c)}+
 \frac1{(b-a)(b-c)}+
 \frac1{(c-a)(c-b)}=0.
\]
This proves the last identity in \eqref{pm:energy-identities}.
Also, grouping the two terms belonging to each unordered pair gives
\[
 S\mathbin{\cdot}b
 =\sum_{i<j}\frac{b_j-b_i}{x_j-x_i}.
\]
The chain rule, \(\nabla E=-2S\), and
\eqref{pm:root-equation} now give
\eqref{pm:forced-energy-balance}; changing sign proves
\eqref{pm:complete-discriminant-balance}.
\end{proof}

The first term in \eqref{pm:forced-energy-balance} is negative, but the
second has no sign under the stated assumptions.  Positivity of the
pair matrix does not remove it.  In fact, the canonical application in
Corollary~\ref{pm:theta-cumulative-cancellation} shows that the physical
pair energy tends to positive infinity for every fixed late-time block.
Consequently that energy cannot be nonincreasing on an entire late-time
tail of the actual theta motion.  This conclusion needs no differentiation
of an asymptotic location error.

\subsection{Normalized matrix evolution and a centered contraction bound}
\label{pm:normalized-subsection}

Let \(\alpha(t)>0\) be continuously differentiable and put
\(y_i=\alpha x_i\), \(r=\alpha'/\alpha\).  The normalized equation is
\begin{equation}
 y'=2\alpha^2S(y)+f,\qquad f=ry+2\alpha b.
 \label{pm:normalized-root-equation}
\end{equation}
In particular the correctly normalized weight evolution is
\begin{align}
 W_{ij}(y)'
  =-2W_{ij}(y)\biggl[
      r+\frac{4\alpha^2}{(y_i-y_j)^2}
       -2\alpha^2\sum_{k\ne i,j}
          \frac1{(y_i-y_k)(y_j-y_k)}
       +2\alpha\frac{b_i-b_j}{y_i-y_j}\biggr].
 \label{pm:normalized-weight-evolution}
\end{align}
This formula follows either by differentiating \(W_{ij}(y)\) directly
in \eqref{pm:normalized-root-equation}, or by differentiating
\(W_{ij}(y)=\alpha^{-2}W_{ij}(x)\).  In particular the rescaling term
\(r\) is present in the normalized formula, and the physical formula
\eqref{pm:physical-weight-evolution} has no such term.

Fix an ordered reference \(c\in\mathbb R^m\), independent of time.
Let \(P_0v=v-m^{-1}(\sum_i v_i)\mathbf 1\) be orthogonal projection
onto the mean-zero subspace, and define
\begin{equation}
 d=y-c,\quad d_0=P_0d,\qquad
 R=ry+2\alpha b+2\alpha^2S(c).
 \label{pm:centered-residual}
\end{equation}
The residual \(R\) is a complete forced residual.  Treating it as zero,
or estimating only one of its summands, is an additional assertion about
the field and is not part of the next theorem.

\begin{theorem}[Centered forced contraction]
\label{pm:centered-contraction}
As long as \(y\) and \(c\) are strictly increasing,
\begin{equation}
 \frac12\frac{d}{dt}\|d_0\|^2
 =-2\alpha^2\sum_{i<j}
       \frac{(d_i-d_j)^2}{(y_i-y_j)(c_i-c_j)}
   +d_0\mathbin{\cdot}P_0R.
 \label{pm:exact-centered-contraction}
\end{equation}
Suppose in addition that \(c_i=c_0+i-1\) and
\(|y_i-c_i|\le\delta(t)<1/2\), with a continuous nonnegative
allowance \(\delta(t)\).  Set
\begin{equation}
 w_0(t)=\frac1{(m-1)(m-1+2\delta(t))},\qquad
 \lambda(t)=2\alpha(t)^2m w_0(t).
 \label{pm:contraction-rate}
\end{equation}
Then, on every compact time interval on which these bounds hold,
\begin{align}
 \|d_0(t)\|
 &\le \exp\left(-\int_{t_0}^t\lambda(u)\,du\right)
             \|d_0(t_0)\|\notag\\
 &\quad+\int_{t_0}^t
     \exp\left(-\int_s^t\lambda(u)\,du\right)
             \|P_0R(s)\|\,ds.
 \label{pm:contraction-convolution}
\end{align}
The mean has its separate equation
\begin{equation}
 \frac{d}{dt}\left(\frac1m\sum_i y_i\right)
  =r\left(\frac1m\sum_i y_i\right)
     +\frac{2\alpha}{m}\sum_i b_i.
 \label{pm:center-equation}
\end{equation}
\end{theorem}

\begin{proof}
The identity
\[
 \frac1{y_i-y_j}-\frac1{c_i-c_j}
  =-\frac{d_i-d_j}{(y_i-y_j)(c_i-c_j)}
\]
gives
\[
 S_i(y)-S_i(c)
 =-\sum_{j\ne i}
       \frac{d_i-d_j}{(y_i-y_j)(c_i-c_j)}.
\]
Both vectors \(S(y)\) and \(S(c)\) have component sum zero.
Subtract their contribution in \eqref{pm:normalized-root-equation},
project onto the mean-zero subspace, and pair with \(d_0\).  Grouping
unordered pairs proves \eqref{pm:exact-centered-contraction}, since
\(d_{0,i}-d_{0,j}=d_i-d_j\).

For clock \(c\), both factors in each mixed denominator have the same
sign.  Their absolute product is at most
\((m-1)(m-1+2\delta)\).  The pair quadratic form is therefore at
least
\[
 w_0\sum_{i<j}(d_i-d_j)^2=w_0m\|d_0\|^2.
\]
Cauchy--Schwarz in the forcing term gives, at times with
\(\|d_0\|>0\),
\[
 \frac{d}{dt}\|d_0\|
 \le-\lambda\|d_0\|+\|P_0R\|.
\]
At a zero of \(d_0\), the same statement holds for the upper right
Dini derivative, directly from the differentiable vector equation.
Equivalently one may regularize the norm and pass to the limit.
Multiplication by the integrating factor yields
\eqref{pm:contraction-convolution}.  Finally, summing
\eqref{pm:normalized-root-equation} and using \(\sum_iS_i=0\)
proves \eqref{pm:center-equation}.
\end{proof}

This theorem identifies the estimate a proposed source-based relaxation
argument must pay: the centered, complete residual \(P_0R\).  It is not
a proof that this residual has a favorable sign or a sufficiently small
norm for theta.  Recovering \(R\) from already-known root motion is an
exact diagnostic, but supplies no new independent decay mechanism.

\subsection{The unforced polynomial model and its Hermite limit}
\label{pm:hermite-subsection}

The polynomial case supplies a rigorous calibration of what deterministic
pair repulsion does achieve.  It also shows why regular late fibers alone
do not exclude a multiple initial zero.

\begin{theorem}[Deterministic Hermite relaxation]
\label{pm:hermite-relaxation}
Let \(P(t,z)\) be a monic real polynomial of degree \(m\) satisfying
\(P_t=-P_{zz}\).  Suppose all its zeros are simple and real at some
\(t_0>0\).  Then they remain simple and real for every \(t\ge t_0\).
Let \(x_1(t)<\cdots<x_m(t)\), put
\(u_i(t)=x_i(t)/\sqrt{2t}\), and set \(s=\log t\).  On the ordered
chamber define
\begin{equation}
 V(u)=\frac14\sum_i u_i^2-\sum_{i<j}\log(u_j-u_i).
 \label{pm:hermite-potential}
\end{equation}
This potential has a unique minimizer \(h\); its coordinates are the
ordered roots of the monic probabilists' Hermite polynomial
\(\operatorname{He}_m\), characterized by
\[
 Q''(z)-zQ'(z)+mQ(z)=0.
\]
The normalized motion and its estimates are
\begin{align}
 \frac{du}{ds}&=S(u)-\frac12u=-\nabla V(u),
 \label{pm:hermite-gradient-flow}\\
 \frac{d}{ds}V(u(s))&=-\|\nabla V(u(s))\|^2,
 \label{pm:hermite-dissipation}\\
 \|u(t)-h\|&\le\left(\frac{t_0}{t}\right)^{1/2}
                              \|u(t_0)-h\|,
 \label{pm:hermite-distance}\\
 V(u(t))-V(h)&\le\frac{t_0}{t}
                              \bigl(V(u(t_0))-V(h)\bigr).
 \label{pm:hermite-energy-decay}
\end{align}
\end{theorem}

\begin{proof}
Initially the simple roots solve \(x'=2S(x)\), because the full
polynomial has unit \(U=1\).  Differentiating \(x=\sqrt{2t}\,u\)
and changing to \(s=\log t\) gives
\eqref{pm:hermite-gradient-flow}.  Direct differentiation of
\eqref{pm:hermite-potential} gives
\[
 \nabla V=\frac12u-S(u),\qquad
 \nabla^2V=\frac12I+\mathcal L(u)\succeq\frac12I.
\]
The potential tends to positive infinity at a finite collision.
It also tends to positive infinity as \(\|u\|\to\infty\): the
quadratic term dominates the at-most-logarithmic negative contribution
from each of the finitely many pair distances.  Its sublevel sets are
therefore compact subsets of the ordered chamber.  Existence of a
minimizer follows, and the displayed strict convexity on this convex
chamber makes the minimizer unique.

At this minimizer, \(2S_i(h)=h_i\).  If
\(Q(z)=\prod_i(z-h_i)\), the same elementary root calculation used
above yields \(Q''(h_i)/Q'(h_i)=2S_i(h)=h_i\).  Thus
\(Q''-zQ'+mQ\) vanishes at all \(m\) distinct roots of \(Q\).
Its leading term cancels, so its degree is at most \(m-1\); it is
identically zero.  The monic polynomial satisfying this equation is
unique by comparing coefficients in descending order.  It is precisely
\(\operatorname{He}_m\).  This also proves its real, simple root
property here, rather than importing that property into the minimization
argument.

The chain rule gives \eqref{pm:hermite-dissipation}.  The normalized
trajectory cannot leave its initial compact sublevel set, so the smooth
gradient equation extends for every finite forward \(s\), with all
gaps positive.  Its coordinates reconstruct simple real zeros of the
original polynomial: both obey the same locally unique root equation,
and continuation gives their agreement.  Hence no later finite collision
or loss of reality occurs.

Strong convexity gives
\[
 (u-h)\mathbin{\cdot}(\nabla V(u)-\nabla V(h))
       \ge\frac12\|u-h\|^2.
\]
Differentiate squared distance along the gradient flow and integrate to
obtain \eqref{pm:hermite-distance}.  Also strong convexity at \(u\),
with comparison point \(h\), implies
\[
 V(u)-V(h)
 \le \nabla V(u)\mathbin{\cdot}(u-h)-\frac14\|u-h\|^2
 \le\|\nabla V(u)\|^2.
\]
Together with \eqref{pm:hermite-dissipation} this gives exponential
decay of the potential gap in \(s\), which is exactly
\eqref{pm:hermite-energy-decay} in physical time.
\end{proof}

For a direct event calibration, the finite heat series gives
\begin{equation}
 e^{-t\partial_z^2}z^m
 =\sum_{k=0}^{\lfloor m/2\rfloor}
       \frac{(-t)^k m!}{k!(m-2k)!}z^{m-2k}
 =(2t)^{m/2}\operatorname{He}_m\!\left(\frac z{\sqrt{2t}}\right).
 \label{pm:exact-hermite-birth}
\end{equation}
Its normalized configuration is exactly \(h\) for every \(t>0\),
while its time-zero polynomial has one zero of multiplicity \(m\).
Perfect equilibrium of the normalized pair statistic therefore does not
erase a past event.  This is a statement about the explicit polynomial
class, not an exclusion of an arithmetic argument for the canonical
source.  The deterministic minimizer in this theorem is one configuration;
no random-matrix distribution has been produced.

The stochastic comparison is different: Dyson Brownian motion includes
diffusion and, in its stationary Gaussian version, a confining drift.
Its equilibrium is a probability distribution, whereas the gradient
flow above converges to one minimizing configuration; see
Joyner--Smilansky~\cite{sc:JoynerSmilansky2015} for the stochastic equations.
